\documentclass[11pt]{article}

% Current official ACL layout. Leave acl.sty and acl_natbib.bst unchanged
% in the same project directory. 
\usepackage{acl}
\usepackage{times}
\usepackage{latexsym}
\usepackage[T1]{fontenc}
\usepackage[utf8]{inputenc}
\usepackage[english]{babel}
\usepackage{microtype}
\usepackage{inconsolata}
\usepackage{graphicx}
\usepackage{amsmath,amssymb}
\usepackage{booktabs}
\usepackage{xcolor}
% Additions and revised passages relative to paper_eng_v4.tex appear in blue.
\newcommand{\revision}[1]{{\color{blue}#1}}

% Enforce English designations also where acl.sty specifies headings.
\usepackage{etoolbox}
\makeatletter
\AtBeginDocument{%
  \renewcommand{\abstractname}{Abstract}%
  \patchcmd{\thebibliography}
    {\section*{References\@mkboth {References}{References}}}
    {\section*{References\@mkboth {References}{References}}}
    {}{}%
}
\makeatother

\title{Comparison of Text Representations and Chunk Selection with Abstention for Evidence Retrieval from German Examination Regulations}

% TODO before submission: Replace placeholders with the final author information.
\author{Seminar Participant / Author \\
  Institut für Sprache und Information \\
  Heinrich-Heine-Universität Düsseldorf \\
  \texttt{TODO: University Email Address}}

\begin{document}

\maketitle

\begin{abstract}
This paper investigates the retrieval of evidence passages in a hierarchically structured German examination regulation. In addition to answerable questions, the task includes an explicit abstention decision for questions that cannot be answered from the source document. The project-specific dataset contains 720 questions, including 120 deliberately unanswerable zero-gold questions, and 201 retrieval chunks. The sparse TF-IDF representation and two fixed multilingual dense representations, Sentence-BERT and multilingual E5, are compared. Building on these, four families of selection methods are evaluated: top-$k$, absolute threshold, top-$k$ with absolute threshold, and selection via a relative score margin. The hyperparameters are determined exclusively on the development split. Five finalists are subsequently frozen for validation; the winner selected there is evaluated exactly once on the held-out test split. The selected configuration, comprising multilingual E5 and top-1 with threshold ($\theta=0.83959$), achieves a mean question-wise chunk-selection F1 score of $0.647$ on the development split, $0.611$ on the validation split, and $0.491$ on the test split. The complete review of the 74 test cases that were not exactly correct identifies 21 retrievals of similar but incorrect chunks, 18 threshold-induced abstentions for answerable questions, 13 false-positive retrievals for zero-gold questions, 9 structural top-1 limitations, 9 cases involving valid alternative evidence passages, and 4 problems involving question wording or lexis. An evidence group unseen in the development and validation splits comprises eight test questions and results in F1 $=0$ in all eight cases. \revision{This illustrates evaluation on gold evidence sets absent from development and validation; the grouping prevents identical non-empty gold sets from crossing split boundaries. A follow-up experiment removes ten problematic evaluation items and uses a random question-level split. Its validation-selected E5 top-1 relative-margin configuration achieves test Q-F1 $=0.650$. Because dataset membership, splitting, and selection differ, this result supplements rather than replaces the original evaluation.}
\end{abstract}


\section{Introduction}

Examination regulations combine general procedural rules, subject-specific provisions, and tabular study-plan information in a hierarchically structured document. Answering a specific question therefore requires more than finding a thematically similar passage. The selected evidence passage must actually contain the rule in question while also being assigned to the correct document level. Passages that use the same technical terms but serve different functions, such as general requirements and their subject-specific implementation, are particularly challenging. Short table cells and entries from study plans can also often be interpreted unambiguously only in conjunction with their headings.

This paper treats the identification of such evidence passages as a separate retrieval task. The prediction target is a set of chunk IDs that provides minimal and complete evidence for a question; the formulation of a natural-language answer is not evaluated. This distinction allows errors in evidence selection to be measured separately from errors made by a downstream reader or generator. Retrieval can serve as an explicit non-parametric knowledge component in knowledge-intensive QA and RAG systems \citep{lewis2020rag}; dense bi-encoder methods also implement this separation for open question-answering tasks \citep{karpukhin-etal-2020-dense}. However, providing an evidence passage implies neither that an answer generated later is correct nor that hallucinations are prevented. The investigation is limited to whether the evidence required for an answer is found.

A pure ranking task does not fully represent the application scenario. Plausibly worded questions may presuppose facts that are not regulated in the examination regulation. A system that always outputs at least one chunk inevitably produces an unsupported retrieval in such cases. The task therefore also requires the explicit prediction of the empty set $[\,]$. The scientific relevance of such no-answer decisions is established in question answering with deliberately unanswerable questions \citep{rajpurkar-etal-2018-know}. In the present retrieval setting, this idea is applied to evidence sets: an abstention is correct exactly when no gold chunk exists for the question.

This problem gives rise to three research questions:

\begin{enumerate}
    \item[\textbf{RQ1}] How do sparse lexical and dense semantic representations differ in their ability to retrieve answer-supporting chunks from a German examination regulation?
    \item[\textbf{RQ2}] How do different strategies for chunk selection and abstention affect retrieval performance and abstention behavior?
    \item[\textbf{RQ3}] How well does the selected configuration generalize to the held-out test split, and which observed error patterns and properties of the data split and evaluation are associated with the remaining errors?
\end{enumerate}

To answer these questions, TF-IDF, a multilingual Sentence-BERT checkpoint, and multilingual E5 are compared using the same cosine-similarity measure. Four selector families separate the quality of the representation from the rule that determines which chunks are output or discarded. The dataset comprises 720 questions and 201 chunks from the 2026 examination regulation for the bachelor's degree program in computational linguistics. \revision{In the original experiment, questions with the same non-empty gold evidence set remain in the same partition. This prevents overlap of identical gold sets, although distinct sets can still share individual chunks.} The paper thus provides a controlled case study of a single German regulatory document: it compares three representations and four selection principles, evaluates an explicit abstention option, and analyzes all test predictions that are not exactly correct.

\revision{The original experiment remains the primary evaluation. A separately reported follow-up examines a revised question set under random question-level splitting, a setting relevant to new questions about an already indexed document. Its purpose is to contextualize the findings under a different evaluation protocol, not to isolate the effect of removing evidence grouping.}


\section{Motivation and Related Work}

\subsection{Sparse and Dense Text Representations}

The classical vector space model represents documents and queries as vectors whose similarity can be determined, among other methods, using cosine similarity \citep{salton1975vector}. TF-IDF weights terms based on their occurrence in a text and their specificity within the document collection \citep{sparckjones1972statistical}. For an examination regulation, this principle provides a natural basis for comparison because questions and evidence passages often use the same technical terms, module names, or regulatory terminology. However, the experiment does not isolate whether this lexical structure actually causes the measured performance; it can serve only as a plausible explanation for the observed TF-IDF results.

Dense representations map texts into continuous vector spaces and can thus capture similarity even when the query and passage do not use the same surface forms. Sentence-BERT produces sentence representations using pair or triplet structures that can be compared directly through cosine similarity \citep{reimers-gurevych-2019-sentence}. For the multilingual variant, an English-language Sentence-BERT model serves as a reference: a multilingual model is trained on translated sentence pairs so that sentences with the same content receive similar vectors regardless of their language \citep{reimers-gurevych-2020-making}. In the present comparison, Sentence-BERT represents a general multilingual sentence-embedding checkpoint. E5, by contrast, was trained contrastively on text pairs and explicitly developed for tasks involving universal vector representations, including retrieval \citep{wang2022text}. The multilingual E5 family extends this training principle to many languages \citep{wang2024multilingual}. This difference in training objective motivates the comparison but does not in itself provide a causal explanation for a project-specific result.

At the same time, broad benchmarks show why neither sparse nor dense methods should be regarded as fundamentally superior across tasks and corpora. BEIR evaluates zero-shot retrieval across heterogeneous datasets and identifies BM25 as a robust lexical baseline \citep{thakur2021beir}. MTEB broadens the perspective to different embedding tasks, datasets, and languages; no single representation method dominates across all tasks \citep{muennighoff-etal-2023-mteb}. Rather than adding another general benchmark to these broad comparisons, the present study examines in a controlled manner how three specified representations behave in a small, terminologically narrow corpus.

\subsection{Abstention and the Legal Retrieval Context}

Selective prediction gives a prediction system the option to reject cases based on a selection criterion. \citet{pmlr-v97-geifman19a} treat this rejection option as a separate learning problem; the present paper does not train such a model but instead uses a simple selection rule based on similarity scores. In question answering, SQuAD~2.0 shows that unanswerable but superficially plausible questions require a separate decision \citep{rajpurkar-etal-2018-know}. Applied to evidence retrieval, this means that the ranking of existing chunks is not the only relevant factor. The system must also recognize when no chunk should be output as evidence.

Work on German legal information retrieval illustrates that matching queries to legal sources is studied as a distinct evaluation problem. GerDaLIR links search queries to relevant German court decisions and compares lexical and neural methods for this purpose \citep{wrzalik-krechel-2021-gerdalir}. GerLayQA links legal questions phrased in lay terms and lawyers' answers to specific statutory sections \citep{buttner-habernal-2024-answering}. For large collections of legal documents, \citet{reuter-etal-2025-towards} also investigate how the representation of structurally similar documents affects retrieval reliability in RAG systems. These tasks differ substantially from the examination regulation considered here in terms of corpus size, document type, and target variable. What they have in common, however, is that the quality of the basis for an answer can be assessed separately by linking it to a documented legal or regulatory source. The present paper therefore positions itself as a narrowly defined case study of evidence selection and abstention, not as a general study of legal NLP.


\section{Methodology}

\subsection{Text Representations and Similarity}

All representations map the same 720 questions and 201 chunks.

\revision{These dimensions refer to the original experiment. The follow-up reuses the existing similarity matrices, selecting the 710 retained question rows by question ID; the 201 retrieval chunks and representation methods remain unchanged.} For TF-IDF, \texttt{TfidfVectorizer(lowercase=True, ngram\_range=(1,3))} is used. The vocabulary is learned exclusively from the chunk texts; the questions are then transformed into the same feature space. The configuration thus considers unigrams, bigrams, and trigrams and converts characters to lowercase. Term-frequency scaling is not enabled.

The checkpoint \nolinkurl{sentence-transformers/paraphrase-multilingual-MiniLM-L12-v2} serves as a general dense sentence representation. Its 384-dimensional vectors are generated through the \texttt{SentenceTransformer} interface, with no custom pooling method added. The retrieval-oriented alternative is \texttt{intfloat/multilingual-e5-base}, which uses 768 dimensions; in accordance with the intended asymmetric use, questions receive the prefix \texttt{query:} and chunks the prefix \texttt{passage:}. Neither of the two dense models is fine-tuned on the project corpus.

For each representation, cosine similarity
\begin{equation}
s(q,c)=\frac{\mathbf{v}_q^\top\mathbf{v}_c}
{\lVert\mathbf{v}_q\rVert_2\,\lVert\mathbf{v}_c\rVert_2}
\end{equation}
is calculated between every question $q$ and every chunk $c$, producing a matrix with $720\times201$ entries for each representation. For the neural encoders, \texttt{normalize\_embeddings=False} is set; the stored vectors are therefore not L2-normalized in advance, and normalization instead takes place as part of the cosine calculation.

\subsection{Selection Rules}

The similarity matrix determines only a ranking or score distribution. Four selector families determine which chunks are actually predicted. For a given question, pure top-$k$ returns the $k$ chunks with the highest similarity scores:
\begin{equation}
S_{\mathrm{top}\text{-}k}(q)=\operatorname{arg\,top}_{k}
\{s(q,c)\mid c\in\mathcal{C}\}.
\end{equation}
As long as at least $k$ candidates exist, the output is never empty. The selector therefore cannot refuse selection without an additional rule.

An absolute threshold $\theta$, by contrast, selects all chunks whose similarity reaches the cutoff:
\begin{equation}
S_{\theta}(q)=\{c\in\mathcal{C}\mid s(q,c)\geq\theta\}.
\end{equation}
This rule can output the empty set, one chunk, or multiple chunks, thereby enabling abstention without limiting the size of the evidence set.

The combined selector first orders all chunks with $s(q,c)\geq\theta$ by descending similarity and returns at most the first $k$ of them, thus combining an absolute minimum similarity with a fixed upper bound on the output size. For $k=1$, the result is either exactly one evidence passage or $[\,]$.

The relative margin evaluates the separation at the decision boundary rather than the absolute score level. With $s_k$ denoting the similarity score at rank $k$ and $s_{k+1}$ the following score,
\begin{equation}
m_k(q)=\frac{s_k-s_{k+1}}{\max(|s_k|,10^{-12})}.
\end{equation}
The hyperparameter $\gamma$ specifies the minimum required size of this relative distance. If $m_k(q)\geq\gamma$, the first $k$ chunks are output; otherwise, the prediction is $[\,]$. For $k=1$, the rule measures how clearly the best chunk is separated from the second-best candidate in relative terms. A high absolute peak score is neither necessary nor sufficient for this criterion.

\subsection{Determination of the Search Ranges}

Because the three representations produce different similarity distributions, no shared threshold ranges are used. The initial regions are derived from robust quantiles of three score distributions on the development split: the respective lowest similarity scores of required gold chunks, the highest scores of non-gold chunks, and the highest scores for zero-gold questions. This procedure restricts the search to regions in which correct retrieval and abstention can be distinguished without using validation or test data. The documented initial ranges are $0.8331$ to $0.8865$ for E5, $0.0591$ to $0.5162$ for TF-IDF, and $0.4740$ to $0.7595$ for Sentence-BERT. All evaluations used to select $k$, $\theta$, and $\gamma$ are conducted exclusively on the development split.

\revision{The ranges above describe the original experiment. The follow-up uses a separate development search and forwards one development-selected configuration per representation and selector family to validation; its results are reported separately below.}


\section{Experimental Setup}

\subsection{Corpus, Chunking, and Gold Standard}

The corpus is based on the 2026 examination regulation for the bachelor's degree program in computational linguistics. The active retrieval collection contains 201 chunks totaling 10{\,}194 words, with an average of 50.72 words and a median of 41 words per chunk; the range is 8 to 204 words. In terms of content, 133 chunks belong to the general examination regulation, 30 to the subject-specific appendix, 36 to the exemplary study plan, and 2 to the front matter. Short excerpts from tables or appendices include their hierarchical context headings so that their meaning remains recognizable outside the original page.

The question dataset comprises 720 entries. Of these, 600 can be answered on the basis of the corpus, while 120 zero-gold questions deliberately have no supporting passage. For each answerable question, \texttt{all\_required\_chunk\_ids} denotes a minimally intended, complete evidence set. This target thus contains the chunks jointly required for the intended answer but does not claim to capture every semantically equivalent alternative evidence passage in the document. Table~\ref{tab:gold-distribution} shows the distribution of empty, single-part, and two-part gold sets.

Chunking of the examination regulation was performed manually. Question creation and gold-evidence assignment were supported by a language model, whereas quality control was conducted manually.

\revision{For the follow-up, dataset version \texttt{manual\_review\_v1} excludes ten questions, leaving 710 entries: 590 answerable and 120 zero-gold questions. Q0485, Q0488, Q0491, Q0494, Q0497, Q0503, Q0515, Q0521, and Q0527 ask whether a named module has a module examination; an alternative study-plan chunk independently supports the same answer as the annotated gold chunk. These exclusions address incomplete coverage of valid evidence alternatives, rather than inherently unanswerable questions. Q0529 is excluded because the formal appendix and exemplary study plan give conflicting semester-hour values (SWS), while the question does not identify the intended source. The retrieval chunks are unchanged. The exclusions and their rationale are documented in \nolinkurl{data/manual_review_v1/README.md}; they do not retroactively alter the original scores.}

\begin{table}[t]
\centering
\small
\begin{tabular}{lrrrr}
\toprule
\textbf{Split} & \textbf{Total} & \textbf{0 Gold} & \textbf{1 Gold} & \textbf{2 Gold} \\
\midrule
Development & 432 & 74 & 340 & 18 \\
Validation & 144 & 24 & 116 & 4 \\
Test         & 144 & 22 & 112 & 10 \\
\midrule
Total       & 720 & 120 & 568 & 32 \\
\bottomrule
\end{tabular}
\caption{Distribution of gold evidence sets across the three data partitions.}
\label{tab:gold-distribution}
\end{table}

\subsection{Data Split and Model Selection}

The data are divided into 432 development, 144 validation, and 144 test questions. The governing grouping rule is \texttt{same\_non\_empty\_gold\_chunk\_set}: all questions with the same non-empty gold evidence set remain in the same data partition. This prevents an evidence passage tested through several paraphrased questions in the test split from having already appeared in a development or validation question with the same gold set. The split contains 315 groups, the largest of which comprises 13 questions, and no non-empty gold set occurs in more than one split. Zero-gold questions each form their own group because they would otherwise be combined into a single large group solely on the basis of their identical empty sets.

\revision{This grouping makes the generalization criterion stricter by withholding identical non-empty gold sets from development and validation. All 201 chunks nevertheless remain available in the retrieval collection: an unseen gold set is not an unseen indexed document. The rule does not exclude all overlap between individual chunks belonging to different gold sets. Although the follow-up uses an ungrouped split, it also changes dataset membership and model selection, so the independent effect of grouping cannot be quantified from the two experiments.}

Selector families and hyperparameters are compared exclusively on the development split. The primary ranking criterion is the mean question-wise F1 score; ties are broken by exact match, mean question-wise precision, and, finally, a lower average number of selected chunks. Five finalists are frozen before validation. Only these configurations are compared on the validation split; no readjustment based on validation takes place. The winner selected there is then run unchanged exactly once on the test split.

\revision{The follow-up uses a seeded random split at the individual-question level (seed 42, proportions 60/20/20), with 426 development, 142 validation, and 142 test questions. No evidence grouping or answer-level stratification is applied. The respective counts of zero-, one-, and two-chunk gold sets are 78/330/18, 22/112/8, and 20/116/6. Twelve finalists, one per representation and selector family, are selected on development and compared on validation. Thus, unlike the original five-finalist comparison, the follow-up validation includes all three representations and all four selector families. The validation-selected winner is frozen before its test evaluation. The selection is recorded in \nolinkurl{frozen_winner.json} and \nolinkurl{validation_ranking.csv} under \nolinkurl{artifacts/experiments/manual_review_v1/outputs/validation/finalists_full_search_v2/}.}

\subsection{Metrics and Comparison Methods}

For each question, the predicted set $P_q$ and gold set $G_q$ are compared. If both sets are non-empty, precision and recall are given by $|P_q\cap G_q|/|P_q|$ and $|P_q\cap G_q|/|G_q|$, respectively; their harmonic mean is the question-wise F1 score. For $P_q=G_q=[\,]$, precision, recall, and F1 are assigned a value of 1 because the system abstains correctly. If exactly one of the two sets is empty, F1 is instead 0. The primary metric, Q-F1, is the arithmetic mean of these question-wise scores. It thus weights each question equally, regardless of the size of its gold set.

A completely missing prediction record must be semantically distinguished from an explicit empty prediction. Missing entries are excluded from the means and counted separately as unanswered. For the frozen test evaluation, prediction records are available for all 144 questions. In addition, mean question-wise precision and recall, exact match (EM), micro-F1, zero-gold abstention rate, rate of empty selections, and average number of selected chunks are reported.

Three simple comparison methods contextualize the test performance. \emph{Random top-1} selects one uniformly random chunk per question; the mean and standard deviation across the fixed seeds 0 to 99 are reported. The \emph{most frequent development answer} predicts the most frequent non-empty exact gold set in the development split for every validation and test question. \emph{Always no selection} returns $[\,]$ regardless of the question. This final method directly measures the Q-F1 attainable solely from the proportion of zero-gold questions.

The experiments were conducted using Python~3.11. Packages central to reproduction include \texttt{scikit-learn==1.9.0}, \texttt{sentence-transformers==3.4.1}, \texttt{transformers==4.46.3}, \texttt{torch==2.1.2} as a CPU build, and \texttt{numpy==1.26.4}. The complete frozen dependency list is documented in the project artifact \texttt{requirements.txt}.


\section{Results and Evaluation}

\revision{Sections~5.1--5.3 report the original grouped experiment, including its baselines and manual error audit. Section~\ref{sec:follow-up} reports the revised-dataset experiment separately; its scores do not replace the original results.}

\subsection{Development Results}

Table~\ref{tab:development-results} reports the best configuration found on the development split for each combination of representation and selector family, with E5 consistently achieving the highest Q-F1 scores among the evaluated representations. The best E5 configuration combines top-1 with the absolute threshold $\theta=0.83959$ and achieves Q-F1 $=0.6474$. The difference relative to pure top-1, with Q-F1 $=0.6165$, is $0.0309$; relative to the best pure threshold selector, with Q-F1 $=0.4941$, it is $0.1533$. The relative margin lies between these configurations, with Q-F1 $=0.6312$.

TF-IDF remains competitive despite the absence of dense semantics: its best configuration is the relative margin, with Q-F1 $=0.5702$, narrowly ahead of top-1 with threshold at $0.5671$. Thus, the ordering of the selectors is not consistent across representations. The evaluated Sentence-BERT checkpoint achieves its best Q-F1 of $0.3989$ with top-1 and threshold and ranks behind E5 and TF-IDF throughout the experimental setup. This result applies only to the specific non-fine-tuned checkpoint and the present dataset.

The differences between the selectors show the trade-off between retrieval and abstention. Pure top-1 effectively limits the output but cannot abstain on any question. Pure threshold selection enables empty and multi-part outputs but exhibits a marked gap between precision and recall for all three representations. For E5, for example, a mean precision of $0.4578$ contrasts with a recall of $0.6285$. Top-1 with threshold combines the abstention option with an upper bound of one chunk; the relative margin, by contrast, uses the local separation between the two highest scores. In the development results, which rule is more favorable depends on the representation.

\begin{table*}[t]
\centering
\small
\begin{tabular}{llcrrrr}
\toprule
\textbf{Representation} & \textbf{Selector} & \textbf{Parameter} & \textbf{Q-F1} & \textbf{EM} & \textbf{Precision} & \textbf{Recall} \\
\midrule
\textbf{E5} & \textbf{Top-1 + Threshold} & $\theta=0.83959$ & \textbf{0.6474} & \textbf{0.6366} & \textbf{0.6528} & \textbf{0.6447} \\
E5 & Relative Margin & $\gamma=0.0025$ & 0.6312 & 0.6204 & 0.6366 & 0.6285 \\
E5 & Top-1 & $k=1$ & 0.6165 & 0.6042 & 0.6227 & 0.6134 \\
E5 & Threshold & $\theta=0.866475$ & 0.4941 & 0.3657 & 0.4578 & 0.6285 \\
\midrule
TF-IDF & Relative Margin & $\gamma=0.235$ & 0.5702 & 0.5579 & 0.5764 & 0.5671 \\
TF-IDF & Top-1 + Threshold & $\theta=0.0999$ & 0.5671 & 0.5532 & 0.5741 & 0.5637 \\
TF-IDF & Top-1 & $k=1$ & 0.5509 & 0.5370 & 0.5579 & 0.5475 \\
TF-IDF & Threshold & $\theta=0.127665$ & 0.4866 & 0.3657 & 0.4596 & 0.5706 \\
\midrule
Sentence-BERT & Top-1 + Threshold & $\theta=0.6751$ & 0.3989 & 0.3958 & 0.4005 & 0.3981 \\
Sentence-BERT & Relative Margin & $\gamma=0.005$ & 0.3866 & 0.3819 & 0.3889 & 0.3854 \\
Sentence-BERT & Top-1 & $k=1$ & 0.3742 & 0.3681 & 0.3773 & 0.3727 \\
Sentence-BERT & Threshold & $\theta=0.7095375$ & 0.3583 & 0.2801 & 0.3343 & 0.4421 \\
\bottomrule
\end{tabular}
\caption{Best development configuration for each representation and selector family. Precision and recall are mean question-wise metrics.}
\label{tab:development-results}
\end{table*}

\subsection{Validation and Held-Out Test}

Among the five finalists frozen in advance, E5 with top-1 and $\theta=0.83959$ achieves the highest Q-F1 on the validation split and is therefore selected as the final system. Sentence-BERT is not among these five finalists; validation therefore does not constitute another complete comparison of all three representations. As Table~\ref{tab:final-system} shows, the selected configuration achieves Q-F1 $=0.6111$, EM $=0.5972$, and a zero-gold abstention rate of $0.4167$ on validation.

On the held-out test split, Q-F1 decreases to $0.4907$, an absolute decline of $0.1204$ relative to validation; EM is $0.4861$, and micro-F1 is $0.5041$. The configuration outputs an average of $0.7917$ chunks and returns the empty set for $20.83\%$ of questions. Of the 22 zero-gold questions, 9 are correctly rejected, corresponding to an abstention rate of $0.4091$. However, the remaining errors cannot be explained by a one-sided tendency to produce either too many or too few outputs: in addition to false-positive retrievals for zero-gold questions, empty predictions occur for answerable questions. Because no repeated splits or comparable uncertainty analysis are available, the difference between validation and test is not interpreted as statistically expected variation.

\begin{table*}[t]
\centering
\small
\begin{tabular}{lrrrrrr}
\toprule
\textbf{Split} & \textbf{Q-F1} & \textbf{EM} & \textbf{Precision} & \textbf{Recall} & \textbf{Micro-F1} & \textbf{ZG Abst.} \\
\midrule
Validation & 0.6111 & 0.5972 & 0.6181 & 0.6076 & 0.6475 & 0.4167 \\
Test         & 0.4907 & 0.4861 & 0.4931 & 0.4896 & 0.5041 & 0.4091 \\
\bottomrule
\end{tabular}
\caption{Frozen winner E5 top-1 with $\theta=0.83959$ on validation and test. For the test, 144 out of 144 predictions are available; TP${}=62$, FP${}=52$, and FN${}=70$.}
\label{tab:final-system}
\end{table*}

The comparison methods in Table~\ref{tab:baselines} show that the test score is not achieved by trivial strategies. Random top-1 achieves a mean Q-F1 of only $0.0047$ across 100 seeds. The most frequent non-empty development answer matches no complete test gold set and therefore achieves Q-F1 $=0$. By contrast, always no selection achieves $0.1528$ because this method handles exactly the 22 zero-gold questions correctly and all 122 answerable questions incorrectly. The final system clearly exceeds these reference points, although its absolute test performance remains limited.

\begin{table*}[t]
\centering
\small
\begin{tabular}{lrrrrr}
\toprule
\textbf{System} & \textbf{Q-F1} & \textbf{Precision} & \textbf{Recall} & \textbf{EM} & \textbf{ZG Abst.} \\
\midrule
Random Top-1 & $0.0047\pm0.0049$ & $0.0050\pm0.0050$ & $0.0045\pm0.0048$ & $0.0041\pm0.0048$ & 0.0000 \\
Most Frequent Development Answer & 0.0000 & 0.0000 & 0.0000 & 0.0000 & 0.0000 \\
Always No Selection & 0.1528 & 0.1528 & 0.1528 & 0.1528 & 1.0000 \\
\textbf{E5 Top-1 + Threshold} & \textbf{0.4907} & \textbf{0.4931} & \textbf{0.4896} & \textbf{0.4861} & \textbf{0.4091} \\
\bottomrule
\end{tabular}
\caption{Test performance of the final system and the frozen comparison methods. For random top-1, the mean and standard deviation across 100 fixed seeds are reported.}
\label{tab:baselines}
\end{table*}

\subsection{Error Analysis and Contextualization}

The manual review covers all 74 test questions without an exact match. Table~\ref{tab:error-analysis} assigns a primary error category to each case. The largest group, comprising 21 cases, involves the retrieval of a similar chunk that is nevertheless incorrect for the specific question. Such confusions correspond to the core problem of the corpus: lexically or thematically similar passages can serve different regulatory functions. A further 18 cases arise because the best correct candidate does not exceed the absolute threshold, causing the system to abstain on an answerable question. Conversely, it retrieves a chunk for 13 unanswerable questions. The similar sizes of these two groups suggest that the errors cannot be fully resolved merely by shifting the threshold globally.

\begin{table}[t]
\centering
\small
\begin{tabular}{lrr}
\toprule
\textbf{Error Category} & \textbf{$n$} & \textbf{Share} \\
\midrule
Similar, Incorrect Chunk & 21 & 28.4\% \\
False Abstention & 18 & 24.3\% \\
False Positive for Zero-Gold & 13 & 17.6\% \\
Top-1 Limitation & 9 & 12.2\% \\
Valid Alternative Evidence Passage & 9 & 12.2\% \\
Question Wording / Lexis & 4 & 5.4\% \\
\midrule
Total & 74 & 100.0\% \\
\bottomrule
\end{tabular}
\caption{Primary categorization of all test cases that are not exactly correct.}
\label{tab:error-analysis}
\end{table}

A second source of error is the fixed top-1 restriction. According to the gold standard, ten test questions require two chunks, so a selector that can output at most one chunk can never achieve an exact match in these cases. If exactly one of the two required chunks is retrieved correctly, this yields $P=1$, $R=0.5$, and thus at most
\begin{equation}
\mathrm{F1}=\frac{2\cdot1\cdot0.5}{1+0.5}=\frac{2}{3}.
\end{equation}
Nine cases that are not exactly correct are assigned primarily to this structural limitation in the audit. The finding thus concerns not only the quality of the E5 representation or the position of the threshold but also the predetermined maximum output size.

\revision{Evidence group G0040 illustrates a cluster of errors under the original grouped protocol, without establishing that grouping caused these errors.} It comprises the eight test questions Q0153 to Q0160 with the shared gold chunk \texttt{PO25CL-GEN-C03-P12-A02}; this non-empty gold set occurs in neither development nor validation. All eight questions achieve F1 $=0$. Q0153, Q0154, Q0155, Q0156, Q0158, and Q0159 each return an incorrect non-empty chunk, while Q0157 and Q0160 abstain with $[\,]$. The cluster is therefore a verified example of failure to generalize to an unseen evidence group. The fact that this evidence group remains entirely unseen is consistent with the deliberately stricter split criterion and does not indicate that the test split was constructed incorrectly.

At the same time, the evaluation is limited by the exhaustiveness of the gold standard. Following the completed review, nine formally incorrect chunk IDs contain a semantically complete alternative evidence passage. These cases are classified as \texttt{valid\_alternative\_evidence} but remain incorrect in the official evaluation because they do not belong to the minimally intended gold set. The official Q-F1 of $0.4907$ is therefore not subsequently replaced. \revision{These observations motivate improvements to the evidence annotations and further evaluation. Difficult evidence groups remain part of the original test result; the follow-up below is reported as a separate experiment rather than as a corrected score.}

For RQ1, the development results thus show a clear project-specific ordering: E5 ranks ahead of TF-IDF, while the evaluated Sentence-BERT checkpoint performs less well. The retrieval-oriented contrastive pretraining of E5 is compatible with this finding but is not tested as an isolated cause. Likewise, the terminological specificity of the examination regulation is a plausible explanation for the relatively strong TF-IDF performance, although the experiment does not provide corresponding causal evidence. For RQ2, there is no universally best selector. Top-1 with absolute threshold is the best E5 configuration and wins on validation; for TF-IDF, the relative margin is slightly ahead on development. The test answers RQ3 by showing weaker generalization than on validation and several types of error operating simultaneously. In addition to ranking and threshold errors, the fixed top-1 limit, unseen evidence groups, and non-exhaustive gold alternatives contribute to the measured result.


\subsection{\textcolor{blue}{Follow-up Evaluation on the Revised Dataset}}
\label{sec:follow-up}

\revision{On the revised development split, E5 again achieves the highest Q-F1 within each selector family. The best development configuration is E5 top-1 with absolute threshold $\theta=0.8396$, with Q-F1 $=0.6487$. Validation instead selects E5 top-1 with relative margin $\gamma=0.0032$. This configuration obtains validation Q-F1 $=0.5540$, compared with $0.5399$ for E5 top-1 with threshold and $0.5329$ for pure E5 top-1. The margin selector's advantage over the threshold finalist is $0.0141$ Q-F1. It determines the winner under the prescribed selection rule but does not establish a generally superior selector. Table~\ref{tab:follow-up} reports the same frozen relative-margin configuration on all three splits; its development score is therefore distinct from the development maximum.}

\begin{table*}[t]
\centering
\color{blue}
\small
\begin{tabular}{lrrrrrr}
\toprule
\textbf{Split} & \textbf{Questions} & \textbf{Q-F1} & \textbf{EM} & \textbf{Precision} & \textbf{Recall} & \textbf{Avg. chunks} \\
\midrule
Development & 426 & 0.6182 & 0.6056 & 0.6244 & 0.6150 & 0.8192 \\
Validation & 142 & 0.5540 & 0.5493 & 0.5563 & 0.5528 & 0.7958 \\
Test & 142 & 0.6502 & 0.6408 & 0.6549 & 0.6479 & 0.7606 \\
\bottomrule
\end{tabular}
\caption{\textcolor{blue}{Revised-dataset results for the validation-selected E5 relative-margin configuration ($k=1$, $\gamma=0.0032$). Precision and recall are question-wise means. All questions have prediction records.}}
\label{tab:follow-up}
\end{table*}

\revision{Test Q-F1 exceeds validation Q-F1 by $0.0962$, or 9.62 percentage points. Different sample composition is a possible explanation, but a single split does not establish the cause or the statistical significance of this difference. The test contains six two-chunk questions, compared with eight on validation, and 20 zero-gold questions, compared with 22 on validation. On test, the winner correctly abstains on 10 of the 20 zero-gold questions and retrieves a chunk for the other ten. It produces 91 exact matches and 51 non-exact predictions. The 74-case manual error taxonomy from the original experiment is not an audit of these new predictions.}

\revision{The random split evaluates new question formulations about the same indexed corpus while allowing evidence targets to recur across partitions. Of the 122 answerable test questions, 110 share their exact non-empty gold evidence set with a development or validation question in the revised split. Because the encoders are fixed, this overlap does not constitute fine-tuning on the test answers. Nevertheless, related questions can influence selector calibration and model selection, so this protocol provides weaker evidence about transfer to gold sets absent from development and validation than the original grouping rule.}

\revision{The test partition is held out within the revised selection workflow, but it is not a fresh independent sample relative to the entire experimental history. Of its 142 questions, 94 previously belonged to development, 28 to validation, and 20 to the original test split, whose errors had already been reviewed. Resplitting therefore does not remove prior exposure. The result is reported as a follow-up evaluation with this qualification, rather than as an independent confirmation on previously untouched questions.}

\revision{The original and follow-up test scores, $0.4907$ and $0.6502$, cannot be interpreted as an isolated improvement due to dataset cleaning or removal of grouping. Question membership, partition assignments, development search, the number of validation finalists, and the selected decision rule differ. The consistent finding is E5's advantage in the development comparisons and its selection on validation in both experiments. The preferred selector and absolute scores are less stable. A controlled comparison would keep the revised question set and search procedure fixed while varying the splitting protocol; repeated prespecified splits could then assess sensitivity to sample composition.}


\section{Conclusion}

\revision{The conclusions below retain the original grouped evaluation as their primary reference; the follow-up supplies additional evidence under the revised protocol.}

The investigation deliberately separates the selection of evidence passages from subsequent answer generation and compares three text representations under the same retrieval and evaluation conditions. With regard to RQ1, multilingual E5 achieves the highest scores among all evaluated representations on the development split; the configuration selected from it also ranks first among the five finalists on validation. Nevertheless, TF-IDF achieves substantially higher scores in the small, terminologically specific corpus than the investigated multilingual Sentence-BERT checkpoint. This ordering applies to the present dataset and the checkpoints used, not to other legal or administrative corpora.

For RQ2, the selection rule proves to be a distinct component of the system. Pure top-1 already achieves strong development results but has no abstention option, whereas a pure threshold can abstain but does not limit the number of output chunks. The combination of E5, top-1, and $\theta=0.83959$ selected on validation combines ranking, an upper bound of one chunk, and the possibility of an empty output. The fact that the relative margin performs slightly better for TF-IDF instead argues against the general dominance of this selector.

RQ3 shows the limitations of the selected configuration. Q-F1 decreases from $0.6111$ on validation to $0.4907$ on the held-out test. The complete review of all 74 cases that are not exactly correct does not attribute this decline to a single cause. Incorrect similar chunks, false abstentions, retrievals for zero-gold questions, the top-1 restriction, alternative valid evidence, and problems with individual question wording occur alongside one another. The eight entirely incorrect questions in the unseen evidence group G0040 illustrate the stricter generalization resulting from the separation of identical gold sets. Evidence retrieval thus remains a clearly measurable upstream component for subsequent QA or RAG systems; the quality of generated answers cannot be inferred from the present results.

\revision{The revised-dataset follow-up also selects E5, but favors top-1 relative margin ($\gamma=0.0032$) rather than the absolute-threshold selector. Its validation and test Q-F1 scores are $0.5540$ and $0.6502$. Together, the experiments support E5's relative strength in this case study while showing that the winning selector and measured performance depend on the evaluation setup. The higher follow-up test score neither supersedes the original result nor identifies which experimental change accounts for the difference.}


\section{Limitations and Future Work}

\revision{The original error counts and abstention figures discussed below refer to the grouped experiment. The follow-up has additional limitations from the reuse of previously evaluated questions and the overlap of evidence targets across its partitions. Its test set is held out only within the revised workflow, and the simultaneous changes in dataset construction and selection preclude causal attribution of the score difference. Neither experiment includes repeated splits or an uncertainty analysis; future comparisons should prespecify the protocols and report variation across splits, with dependence between questions sharing evidence taken into account.}

The empirical scope is limited to one regulatory document for one degree program. Transferability to other examination regulations or universities, as well as to larger collections of legal documents, was not investigated. The questions and gold evidence sets were constructed specifically for the project with language-model support and manually reviewed by the author. Because no independent second annotation was conducted, inter-annotator agreement could not be assessed.

The evaluation target represents a minimally intended complete evidence set but not an exhaustive set of all permissible evidence passages.

\revision{Removing nine questions with valid alternative passages in the follow-up reduces one known annotation problem but does not establish that all remaining gold sets exhaust the admissible evidence. A future annotation scheme should represent alternative sufficient evidence sets explicitly, allowing such questions to remain in the benchmark.} The nine test predictions reviewed as valid alternatives show the practical consequence of this decision. They do not justify a subsequent change to the official test score, but they support future gold standards that explicitly model equivalent evidence sets or canonicalized references. In addition, the final top-1 selector limits the attainable performance for questions with multiple required evidence passages; this concerns ten questions in the test split alone.

Abstention is also not conclusively calibrated. The system identifies 9 of 22 zero-gold questions but incorrectly retrieves a chunk for the remaining 13. At the same time, 18 threshold-induced abstentions occur for answerable questions. Characterizing the system merely as too conservative or too permissive therefore falls short. Future experiments should investigate calibration separately by question type and permit selectors that output zero, one, or multiple chunks depending on the question.

The frozen comparison includes neither BM25 nor hybrid sparse-dense methods, cross-encoder reranking, or generative reader models. BEIR identifies BM25 as a robust zero-shot baseline \citep{thakur2021beir}; however, a direct comparison on the present corpus remains outstanding. Meaningful extensions therefore include BM25 and hybrid methods, downstream reranking of similarly scored chunks, adaptive evidence sets, and content-type-aware strategies for running text, appendices, and tables. Downstream answer generation should be evaluated only after such an extension of retrieval. RAG provides a suitable system framework for this purpose \citep{lewis2020rag}, but answer correctness and faithfulness to evidence would then have to be measured in addition to, and separately from, the retrieval performance investigated here.

\bibliography{references}

\end{document}
